% Schematic of the physical contact set-up, drawn from photographs of the rig.
% (a) side elevation of the arrangement, (b) the contact geometry it defines.
% Uses only tikz + arrows.meta/calc, which config/packages.tex already loads.
\begin{tikzpicture}[
    font=\small,
    wood/.style={draw=brown!70!black, fill=white},
    woodDark/.style={draw=brown!70!black, fill=white},
    metal/.style={draw=gray!70!black, fill=white},
    robot/.style={draw=gray!60, fill=white},
    frameArrow/.style={-{Latex[length=2mm]}, thick},
  ]

% ============================ (a) side elevation =========================
\begin{scope}

  % ---- last robot link and flange -------------------------------------
  \draw[robot, rounded corners=6pt] (-0.75,5.75) rectangle (0.75,7.30);
  \node[gray!55!black] at (0,6.55) {\footnotesize link 7};
  \draw[metal] (-0.60,5.42) rectangle (0.60,5.75);

  % ---- gripper body ---------------------------------------------------
  \draw[robot, rounded corners=3pt] (-0.95,4.52) rectangle (0.95,5.42);
  \node[gray!55!black] at (0,4.97) {\footnotesize hand};

  % ---- custom aluminium fingers ---------------------------------------
  \draw[metal] (-0.62,4.52) -- (-0.74,3.74) -- (-0.44,3.74) -- (-0.34,4.52) -- cycle;
  \draw[metal] ( 0.62,4.52) -- ( 0.74,3.74) -- ( 0.44,3.74) -- ( 0.34,4.52) -- cycle;
  \node[gray!50!black, anchor=west] at (0.86,4.12) {\footnotesize custom fingers};

  % ---- wooden tool -----------------------------------------------------
  % Anchored so its lower-LEFT corner sits exactly on the workpiece: the
  % residual orientation error leaves the face a few degrees off parallel,
  % so one edge makes contact first.
  \begin{scope}[shift={(-0.62,2.98)}, rotate=9]
    \draw[woodDark] (0,0) rectangle (1.45,0.42);
    \node[font=\footnotesize] at (0.72,0.21) {tool};
  \end{scope}

  % ---- workpiece and bench --------------------------------------------
  \draw[wood] (-1.75,2.58) rectangle (1.75,2.98);        % raised strip
  \draw[wood] (-2.60,2.18) rectangle (2.60,2.58);        % base board
  \draw[fill=white, draw=gray!50] (-3.20,1.72) rectangle (3.20,2.18);
  \node[anchor=west] at (1.85,2.78) {\footnotesize workpiece};
  \node[anchor=west] at (2.70,2.38) {\footnotesize base board};
  \node[gray!55!black, anchor=west] at (3.30,1.95) {\footnotesize bench};

  % ---- contact marker ---------------------------------------------------
  \fill[red!75!black] (-0.62,2.98) circle (1.7pt);
  \draw[red!70!black, thin] (-0.66,3.02) -- (-1.95,3.60);
  \node[red!60!black, anchor=east, font=\footnotesize] at (-1.95,3.60)
        {leading edge};

  \node at (0,1.15) {(a) side elevation};
\end{scope}

% ============================ (b) contact geometry =======================
\begin{scope}[xshift=9.4cm, yshift=0.15cm]

  % ---- the physical plane, tilted by b ------------------------------------
  \begin{scope}[shift={(0,2.30)}, rotate=-15]
    \draw[wood] (-2.10,-0.24) rectangle (2.10,0);
  \end{scope}

  % ---- horizontal reference and the tilt angle ----------------------------
  \draw[gray!55, dashed] (-2.15,2.30) -- (2.60,2.30);
  \draw[{Latex[length=1.5mm]}-, gray!60!black] (1.85,2.30) arc (0:-15:1.85);
  \node[gray!50!black, font=\scriptsize] at (2.20,2.06) {$b$};
  \node[gray!50!black, font=\scriptsize, anchor=west] at (1.95,2.44) {horizontal};

  % ---- contact point -------------------------------------------------------
  \coordinate (C) at (-1.05,2.58);
  \fill[red!75!black] (C) circle (1.8pt);
  \draw[red!70!black, thin] (C) -- (-2.00,3.35);
  \node[red!60!black, anchor=south east, font=\footnotesize] at (-1.95,3.32)
        {contact};

  % ---- surface frame at the contact ---------------------------------------
  % t_1 points down-slope along the plane; n_s is its normal. Kept short so
  % neither collides with the tool face or the lever arrow.
  \begin{scope}[rotate around={-15:(C)}]
    \draw[frameArrow] (C) -- ++(0.95,0)
          node[anchor=north west, xshift=-1mm, yshift=1mm] {$t_1$};
    \draw[frameArrow] (C) -- ++(0,1.55) node[anchor=south east] {$n_s$};
  \end{scope}

  % ---- the tool face, a few degrees off the plane --------------------------
  \begin{scope}[shift={(C)}, rotate=-4]
    \draw[woodDark] (0,0.02) rectangle (1.25,0.34);
    \node[font=\scriptsize] at (0.35,0.16) {tool};
  \end{scope}

  % ---- configured normal ---------------------------------------------------
  \draw[frameArrow, gray!55!black, densely dashed] (0.05,2.58) -- ++(0.11,1.35);
  \node[gray!50!black, font=\scriptsize, anchor=south] at (0.16,3.95)
        {configured $n$};

  % ---- the TCP, the compliance centre, and the lever between them ----------
  % r_c runs from the TCP to p_c. It starts at neither the contact point nor
  % the selected tool point: those are reached from the TCP by the separate
  % tool-geometry offset, and only the TCP defines this lever. p_c is placed
  % clear of the plane and of the tool face, since the measurements say its
  % height is immaterial and the drawing is free to put it where it reads.
  \coordinate (TCP) at (-0.403,2.876);
  \coordinate (P) at (1.55,3.28);
  \fill (TCP) circle (1.7pt);
  \draw[thin] (TCP) -- (-0.38,3.07);
  \node[font=\scriptsize, anchor=south] at (-0.35,3.10) {$p_{\mathrm{TCP}}$};
  \draw[blue!65!black, thin, -{Latex[length=1.6mm]}] (TCP) -- (P);
  \node[blue!55!black, font=\scriptsize, anchor=north west] at (0.30,2.92) {$r_c$};
  \fill[blue!70!black] (P) circle (1.9pt);
  \node[blue!60!black, anchor=west, font=\footnotesize, align=left]
        at (1.70,3.28) {compliance\\centre $p_c$};

  % Both sub-labels sit on one baseline at absolute y = 1.15. This scope
  % carries yshift=0.15cm, so the node is placed at 1.00 to compensate; if
  % that yshift changes, this value changes with it.
  \node at (0.1,1.00) {(b) contact geometry};
\end{scope}

\end{tikzpicture}
